\documentclass[10pt]{article}
\usepackage[letterpaper,margin=.38in]{geometry}
\usepackage{amsmath,amssymb,multicol}
\pdfinfoomitdate=1
\pdftrailerid{}
\pagestyle{empty}
\setlength{\parindent}{0pt}
\setlength{\parskip}{1pt}
\setlength{\jot}{0pt}
\setlength{\columnsep}{12pt}
\setlength{\abovedisplayskip}{1pt}
\setlength{\belowdisplayskip}{1pt}
\let\bigsum\sum
\renewcommand{\sum}{\bigsum\nolimits}
\newcommand{\fl}[1]{\left\lfloor #1\right\rfloor}
\newcommand{\rng}[2]{\{#1,\ldots,#2\}}
\begin{document}
\noindent\parbox{\textwidth}{\centering
{\large\bfseries A completely specified nine-positive consistency equation}\\[-1pt]
{\small Nine positive inputs; accepted evaluator: 433 operations, degree $1{,}628{,}768$. All fixed data below.}}\par
\begin{multicols}{2}
\fontsize{9}{9.0}\selectfont
The only unknowns are $b,g,h,s,w,C,i,j,n\in\mathbb Z_{>0}$.
All other letters are polynomial abbreviations or fixed constants.
The equation is
\[
\boxed{\bigl[\Theta^2+\Omega^2(a_0-a_1)\bigr]^2
       -4\Omega^2a_0\Theta^2=0.}
\]
\textbf{Source and ratio polynomials.} For any displayed finite list put
$p(e_1,\ldots,e_t)=\sum_{\nu=1}^{t}B^{e_\nu}$. Define
\[
\begin{gathered}
B=2^{864026}b^2,\quad N_0=B^{117},\quad\ell=p(\mathcal W),\quad Z=p(\mathcal T),\\
u=\ell+g,\quad G_c=D_2(B)u^2+D_1(B)u+D_0(B),\\
S_c=g+N_0(Z+2G_c),\\
T_c=N_0-1-\ell(b-1)+N_0Z(B-2),\\
L_c=bB^{2200},\quad R=L_c(S_c+T_c+1)-S_c-1,\\
N=2^{6480184}B^{1162},\quad c_*=bB^{116},\\
U=N^2w,\quad Y=N^2s,\quad M=RY,\\
P=2M^2U,\quad k=R+1+h(P-1),\\
A=M(U+1),\quad B_1=2R+1,\quad W_1=bw,\\
J_r=4A-5,\quad Q_r=3W_1C-2(W_1^2-1),\\
m=k^2(4c_*-g)-16c_*(C-kY)^2,\quad\Omega=J_rm,\\
E_0=J_r(n+W_1^2)+Q_r,\quad V_0=Pk,\\
\delta=A^2-1,\quad D=1+\delta C^2,\quad e=iJ_rC^2,\quad d_e=De,\\
F_p=1+(\delta e)d_e,\quad\Lambda=CD+1-A,\quad G_p=A+F_p\Lambda,\\
H=C+F_p[B_1+(j-1)C],\quad T_m=Ad_e,\quad V_1=T_m(G_pH),\\
I=1+(G_pH)^2-H^2,\quad a_0=1+V_0^2-k^2,\quad a_1=DF_pI,\\
\Theta=E_0-\Omega(V_0+V_1).
\end{gathered}
\]
The complete fixed lists, with $p(\mathcal W)$ and $p(\mathcal T)$ meaning
the sum over all entries, are
\[
\begin{aligned}
\mathcal W={}&(1,7,61,36,16,85,20,67,43,11,34,13,96,80,29,50,\\
 &2,109,110,112,116,91,53,55,27,101,92,9,107,37,60),\\
\mathcal T={}&(28,124,325,384,589,630,643,696,846,1073,1082,\\
 &1159,1251,1261,1450,1510,1627,1756,1793,1807,2028,2050).
\end{aligned}
\]
\textbf{All 85 kernel terms.} Equal-coefficient powers are grouped only by
the explicitly defined finite-list notation $p$:
\[
\begin{aligned}
D_0(B)={}&2K_0B^{124}-2p(325,384,1450,1627)-6B^{630},\\
D_1(B)={}&2p(26,33,477,542,570,605,1020,1027)\\
 &+2p(1047,1067,1144,1234,1354,1358,1370)\\
 &+2p(1501,1691,1719,1733,1756,1918,1941)\\
 &-2p(1,493,662,754,982,1698,1796)-8B^{514}-6B^{1023}\\
 &-2K_0B^{63}-2(K_0+1)B^{117}+2\beta(K_0+1)B^{123}\\
 &-2(M_C+M_F)B^{845}-4\beta B^{1260},\\
D_2(B)={}&-\beta p(0,87,661,843,1795)\\
 &-p(29,83,119,135,338,377,393,758,812,845)\\
 &-p(1076,1176,1256,1295,1512,1632,1794,1807,1843)\\
 &-2p(105,398,611,858,1452,1542)+2p(299,326,1326,1521)\\
 &+p(384,400,691,837,1823,1850)-2\beta M_FB^{844}\\
 &-4B^{907}-2\beta^2B^{1259}.
\end{aligned}
\]

\textbf{The complete fixed coefficient tuple.}
The divisions and floors below define constants only; the rational
intermediate defining $K_0$ is an integer.
\[
\begin{gathered}
\rho=2^{89},\quad\mathcal B=\rho^{4831},\quad\beta=\mathcal B-1,\\
\lambda=2^{29},\quad\theta=2^{42},\quad\mathcal Q=2^{40},\\
G_t(x)=\sum_{r=0}^{t-1}x^r,\quad
\mu=510+2^{17}+2^{26}+2^{28}+\lambda(2^{44}-1).
\end{gathered}
\]
\[
\begin{aligned}
\mathcal S={}&\rng0{456}\cup\rng{495}{504}\cup\{507,508\}\\
 &\cup\rng{517}{537}\cup\rng{539}{549}\\
 &\cup\rng{567}{583}\cup\{1203,1807\}.
\end{aligned}
\]
\[
\begin{aligned}
M_C={}&\mathcal B-1-\sum_{r\in\mathcal S\setminus\{1\}}\rho^r-2\rho^{507},\\
M_F={}&\mathcal B-1+4+\mu(\rho^{2409}+\rho^{3015})\\
 &+\rho^{528}G_{10}(\rho)+\rho^{585}G_{11}(\rho)+\rho^{1203}+\rho^{1807}\\
 &+(2^{87}-1)\rho^{4823}G_4(\rho)+(2^{38}-1)\rho^{4827},\\
K_0={}&\rho^{33}+\rho^{68}+\rho^{664}+\rho^{1268}+\rho^{4877}\\
 &+(\rho^{2409}+\rho^{3015})\Phi(1/\rho),
\end{aligned}
\]
\[
\begin{aligned}
\Phi(x)={}&O G_{457}(x)+\lambda\sum_{a=0}^{456}\ell_a x^{r_a}\\
 &-(255\,2^9+2^{27})x^{374}-(255\,2^{18}+2^{27})x^4\\
 &+\lambda\bigl[(x^{539}+2^{21}x^{517})G_{11}(2x)\\
 &\qquad +(2^{11}x^{528}+2^{32}x^{495})G_{10}(2x)\\
 &\qquad +\theta(x^{539}+x^{1203}+x^{1807})\bigr]\\
 &+\bigl[(2^9+\lambda(2+2^{33}))x^{567}\\
 &\qquad +(2^{18}+\lambda(2^{22}+2^{33}))x^{575}\bigr]G_8(2x)\\
 &+(2^{27}+\lambda(2^{11}+2^{32}))x^{583}.
\end{aligned}
\]
\textbf{The ranks and labels.}
The following 457 fixed quadruples are $t_r=(w_r,n_r,e_r,s_r)$.
From the fixed integer $\mathcal A$ below, for $0\le r<430$ put
\[
\begin{gathered}
b_r=\fl{\mathcal A/1024^r}\bmod1024,\quad u_r=b_r\bmod256,\\
v_r=\fl{b_r/256}\bmod2,\quad d_r=\fl{b_r/512},\\
t_{2+r}=((1-d_r)(512+u_r),512+r,d_r(256+u_r),v_r).
\end{gathered}
\]
Here $a\bmod b=a-b\fl{a/b}$. The remaining quadruples are
\[
\begin{array}{c|c|l}
\text{index}&\text{range}&t_r\\\hline
0&& (2,3,3,1)\\
1&& (5,0,0,0)\\
432+r&0\le r<2&(0,942+r,0,942+r)\\
434+r&0\le r<2&(0,r,0,r)\\
436+r&0\le r<4&(0,(r\bmod2)+942\fl{r/2},0,6)\\
440&&(2,2,2,0)\\
441&&(3,4,4,698)\\
442&&(4,4,4,0)\\
443+r&0\le r<3&(1,6,1,2+r)\\
446+3r+t&0\le r,t<3&(r^2,5,t,5)\\
455&&(256,0,0,512)\\
456&&(0,0,512,512)
\end{array}
\]
For these bound indices define
\[
\begin{aligned}
\kappa_r&=2^{30}w_r+2^{20}n_r+2^{10}e_r+s_r,\\
\ell_r&=2w_r+2^{11}n_r+2^{22}e_r+2^{32}s_r,\\
r_0&=0,\quad r_1=1,\\
r_t&=2+\sum_{r=2}^{456}\fl{\frac{\mathcal Q+\kappa_t-\kappa_r-1}{\mathcal Q}}
 \quad(2\le t\le456),\\
O&=1+255\,2^9+255\,2^{18}+2^{27}+\lambda.
\end{aligned}
\]
Subscripts in the constant recipes are bound indices, not extra unknowns.

\textbf{The integer $\mathcal A$.} Digit values are $0,\ldots,61$ in the alphabet
\begingroup\fontsize{8}{8.8}\selectfont\ttfamily
0123456789ABCDEFGHIJKLMNOPQRSTUVWXYZabcdefghijklmnopqrstuvwxyz\par
\endgroup
For the following 722 characters put
$\mathcal A=\sum_{r=0}^{721}\epsilon_r62^{721-r}$, where $\epsilon_r$
is the value of character $r$, starting at zero:
\begingroup
\fontsize{8}{8.8}\selectfont\ttfamily
RCtYkANMQ1gmPqxuKhLmpi1mjCvu243YjyAxv9hbC2UsurRLykLFhuHVBGGUZ5\par
whVShV0FQcVVQSR2GS6JAW1DGYxGKlA4rJavFCW5i63LLSs8xEGvWTbLjSFULV\par
ODe48Kw2ghEaa5Q9PZZVHtki7C45jfVYROOzvVT9DAGjL0BlBSHQU8WPUF3SpC\par
kTfDQF7mUboMEAeyPWARCyClpnxtqgsit318hQgiIdKauxc6memnYNtycB0s7z\par
K8dMgYaLRxNjlcLvJzZ9Z5wdCG3wjd5P6ZVqC2SXVaimqXrTY5dNxbCH2B5VKA\par
Sp16kht0JV65p4VclYsI6zi6nIr7elpUCOX7IHdu6qKiu6Kt5MqlYZjChEzG6y\par
W1WsPdMB7lPL3oOCFqFncCtVEMUounam1fyhduQYoTyTCIEMqcD2BPQfMVaBa9\par
BZHQ8k7rpS3eiM7Rsak5uUsoP1Tf1uhWshndLuN8yz7UdD5QXM6AJ3ovPdWsqO\par
HGBeRaR8kch4R9JVbkAxS2ByaeE2O6czIT0sWtpasWDdZB9FuDQUFYhSvPPPr8\par
nSQmnMkF90CBzWkpeP81azGOWtiT2XTxpaLcR4ecaK1uZyvBleYknk6QTwECl5\par
XJ3nJBoP4sMTGaV6unoGzq9s1KN4KtXrVRuTpi6EuqizW9njcp9Uikq1VNjVtw\par
W1AGLJu3eaWg0EEF7vROqLuNbzEheYnyi4SORWVp\par
\endgroup
\smallskip
\normalfont\fontsize{8}{8.8}\selectfont
Absence of positive integer solutions is equivalent to
$\mathrm{Con}(\mathrm{ZFC})$ by the inherited source/native calibration and
the ordinary saturated Pell argument. This page specifies the equation;
it does not supply an end-to-end formal Lean proof.
\end{multicols}
\end{document}
